\documentclass[12pt]{article}

\usepackage{amsmath}
\usepackage{amssymb}
\usepackage{geometry}
\usepackage{booktabs}
\usepackage{setspace}

\geometry{margin=1in}
\setstretch{1.2}

\title{CSE400 -- Fundamentals of Probability in Computing\\
Lecture 15--17: Joint Random Variables and Joint Distributions}
\author{Aryan Thakkar // Enrolment no.: AU2440248}
\date{March 10, 2026}

\begin{document}

\maketitle

\section{Why Study Joint Random Variables?}

\subsection{Motivation}

In science and real life we are often interested in two (or more) random variables simultaneously.

Examples include:
\begin{itemize}
\item Height and weight of giraffes
\item IQ and birthweight of children
\item Frequency of exercise and rate of heart disease
\item Air pollution level and respiratory illness rate
\item Number of Facebook friends and age of members
\end{itemize}

These examples illustrate situations where multiple random variables are studied together.

\subsection{Example \#1: Smart Transportation}

Two random variables in a smart traffic system.

Question:

Can speed be high when the traffic density is also high?

Possible answers:
\begin{itemize}
\item Yes
\item No
\end{itemize}

Applications in Intelligent Transportation Systems (ITS):
\begin{itemize}
\item Adaptive traffic signals
\item Congestion prediction
\item Autonomous vehicle routing
\end{itemize}

\subsection{Example \#2: Wireless Network Performance}

Random variables:
\[
X = \text{Vehicle Speed}
\]
\[
Y = \text{Traffic Density}
\]

Another setting:
\[
X = \text{Signal-to-Noise Ratio (SNR)}
\]

Reason for studying jointly:

Throughput depends on SNR, and both vary due to:
\begin{itemize}
\item fading
\item interference
\item congestion
\end{itemize}

Question:

If SNR is high, will throughput always be high?

Possible answers:
\begin{itemize}
\item Yes
\item No
\end{itemize}

Applications:
\begin{itemize}
\item 5G scheduling
\item Adaptive modulation
\item Network planning
\end{itemize}

\subsection{Additional Examples}

Example \#3: Weather Prediction

Application: weather forecasting.

Example \#4: Drone Delivery

Application: logistics and UAV systems.

Example \#5: Rolling Two Dice

Let:
\[
X, Y = \text{values obtained on two dice}
\]

Example event:
\[
X + Y = 7 \quad \text{and} \quad X > Y
\]

This requires a joint PMF.

\subsection{Joint Distribution Concept}

When multiple random variables are considered together:

\begin{itemize}
\item They possess a joint distribution.
\item The joint distribution allows computation of probabilities involving both variables.
\item It helps understand the relationship between variables.
\end{itemize}

Simplest situation:

When variables are independent.

When variables are not independent:

Covariance and correlation are used to measure dependence.

\section{Discrete Case -- Joint PMF}

\subsection{Setup}

Suppose $X$ and $Y$ are discrete random variables.

The ordered pair $(X,Y)$ takes values in the product set:

\[
X = \{x_1, x_2, \dots, x_n\}
\]

\[
Y = \{y_1, y_2, \dots, y_m\}
\]

\[
(X,Y) \in \{(x_1,y_1),(x_1,y_2),\dots,(x_n,y_m)\}
\]

\subsection{Definition: Joint PMF}

The joint probability mass function is defined as

\[
p(x_i,y_j) = P(X = x_i, Y = y_j)
\]

This represents the probability that both outcomes occur simultaneously.

\section{Joint PMF Table Representation}

The joint PMF can be organized in a table format.

\begin{center}
\begin{tabular}{c|cccc}
$X \backslash Y$ & $y_1$ & $y_2$ & $\dots$ & $y_m$ \\
\hline
$x_1$ & $p(x_1,y_1)$ & $p(x_1,y_2)$ & $\dots$ & $p(x_1,y_m)$ \\
$x_2$ & $p(x_2,y_1)$ & $p(x_2,y_2)$ & $\dots$ & $p(x_2,y_m)$ \\
$\vdots$ & $\vdots$ & $\vdots$ & $\ddots$ & $\vdots$ \\
$x_n$ & $p(x_n,y_1)$ & $p(x_n,y_2)$ & $\dots$ & $p(x_n,y_m)$ \\
\end{tabular}
\end{center}

\subsection{Key Properties of Joint PMF}

\begin{enumerate}
\item Probability bounds
\[
0 \le p(x_i,y_j) \le 1
\]

\item Total probability equals one
\[
\sum_{i=1}^{n}\sum_{j=1}^{m} p(x_i,y_j) = 1
\]
\end{enumerate}

\section{Example: Rolling Two Dice}

\subsection{Experiment}

Roll two fair dice.

Define

\[
X = \text{value on the first die}
\]

\[
T = \text{total sum of both dice}
\]

\subsection{Outcome Probabilities}

Each outcome $(i,j)$ occurs with probability

\[
\frac{1}{36}
\]

Thus a joint probability table for $(X,T)$ is constructed.

\subsection{Joint Probability Table Observations}

The table entries consist of either

\[
\frac{1}{36}
\]

or

\[
0
\]

Observations:

\begin{itemize}
\item Entries are either $\frac{1}{36}$ or $0$.
\item $X$ and $T$ are not independent.
\end{itemize}

\section{Continuous Case -- Joint PDF}

\subsection{Relationship to Discrete Case}

Continuous case is analogous to the discrete case:

\begin{itemize}
\item Discrete values $\rightarrow$ Continuous intervals
\item Joint PMF $\rightarrow$ Joint PDF
\item Sums $\rightarrow$ Double integrals
\end{itemize}

\subsection{Product Region}

If

\[
X \in [a,b], \quad Y \in [c,d]
\]

then

\[
(X,Y) \in [a,b] \times [c,d]
\]

\subsection{Definition: Joint PDF}

A function $f(x,y)$ is the joint probability density function of $X$ and $Y$ if

\[
P((X,Y) \text{ lies in a small rectangle } dx\,dy)
\approx f(x,y)\,dx\,dy
\]

\subsection{Key Properties}

\begin{enumerate}
\item Non-negativity
\[
f(x,y) \ge 0
\]

\item Normalization
\[
\int_a^b \int_c^d f(x,y)\,dx\,dy = 1
\]
\end{enumerate}

\subsection{Important Note}

The joint PDF may take values greater than $1$, because it represents a density and not a probability.

\section{Continuous Example}

\subsection{Given}

\[
f(x,y) = 4xy
\]

for

\[
0 \le x \le 1, \quad 0 \le y \le 1
\]

Thus $(X,Y)$ lies in the unit square:

\[
[0,1]^2
\]

\subsection{Tasks}

\begin{enumerate}
\item Show that $f(x,y)$ is a valid joint PDF.
\item Visualize the event
\[
A = \{X < 0.5, Y > 0.5\}
\]
\item Compute
\[
P(A)
\]
\end{enumerate}

\section{Example Solution}

\subsection{Step 1: Check Validity (Normalization)}

To verify that $f(x,y)$ is a valid joint PDF, check

\[
\int_0^1 \int_0^1 4xy \, dx \, dy = 1
\]

This confirms the function is a valid joint PDF.

\section{Joint Cumulative Distribution Function (Joint CDF)}

\subsection{Definition}

The joint cumulative distribution function is defined as

\[
F(x,y) = P(X \le x, Y \le y)
\]

\section{Joint CDF -- Continuous Case}

For continuous random variables:

The joint CDF is obtained by integrating the joint PDF over the appropriate region.

\section{Joint CDF -- Discrete Case}

For discrete random variables:

The joint CDF is obtained by summing the joint PMF over all values satisfying

\[
X \le x, \quad Y \le y
\]

\section{Properties of the Joint CDF}

\begin{itemize}
\item The function is non-decreasing in both variables.
\item Limits correspond to probabilities approaching $0$ or $1$ depending on the region.
\item The function represents cumulative probability over the joint domain.
\end{itemize}

\section*{End of Lecture}

Thank You.

\end{document}